\documentclass[conference,a4paper]{IEEEtran}
\IEEEoverridecommandlockouts
% --- PACKAGES ---
% Pass this option before todonotes/TikZ can load xcolor.
\PassOptionsToPackage{table}{xcolor}
\usepackage{cite}
\usepackage{amsmath,amssymb,amsfonts}
\usepackage{algorithmic}
\usepackage{graphicx}
\graphicspath{{image/supplementary/}}
\usepackage{textcomp}
\usepackage[T5]{fontenc}
\usepackage[utf8]{inputenc}
\usepackage{makecell}
\usepackage[disable,colorinlistoftodos,textsize=tiny]{todonotes}
\usepackage{xcolor}
\usepackage{booktabs}
\usepackage[table]{xcolor}
\usepackage{multirow}
\usepackage{caption}
\usepackage{placeins}
\usepackage{subcaption}
\usepackage{tabularx}
\usepackage{soul}
\usepackage{tikz}
\usepackage{adjustbox}
\usepackage{changes}
\usepackage{flushend}
\usepackage{color}
\flushend
\setlength{\marginparwidth}{2cm}
\renewcommand{\thesection}{S\arabic{section}}
\renewcommand{\thesubsection}{\thesection.\arabic{subsection}}
\renewcommand{\thetable}{S\arabic{table}}
\renewcommand{\thefigure}{S\arabic{figure}}
\captionsetup{font=footnotesize,labelfont=bf}
\newcolumntype{Y}{>{\raggedright\arraybackslash}X}
\setlength{\emergencystretch}{1em}
\begin{document}
\twocolumn[{
\begin{center}
{\Large\bfseries Tài liệu bổ sung}\\[4pt]
{\normalsize cho bài báo ``Ultra-Short PPG Dynamics Across the Wakefulness-to-Drowsiness Transition:
A Nonlinear Time-Series Analysis''}
\end{center}
}]

\renewcommand{\tablename}{Bảng}
\renewcommand{\figurename}{Hình}
\section{Kiểm tra tính toàn vẹn dữ liệu và số window được giữ lại}\label{sec:supp_1}
Dữ liệu gồm 20 session ghi nhận với tổng cộng 1,620,038 mẫu. Session 01 có tần số lấy mẫu 50 Hz; các session còn lại khoảng 25 Hz. Không phát hiện giá trị tín hiệu hoặc thời gian bị thiếu, không phải số hay vô hạn; dấu thời gian không trùng lặp và nhãn đều hợp lệ. Có 69 khoảng gián đoạn dài hơn $1.5$ lần khoảng lấy mẫu trung vị. Không nội suy qua các khoảng gián đoạn, và window phân tích không cắt qua chúng.

Bảng~\ref{tab:supp_retained} liệt kê số window còn lại sau bước sàng lọc nhiễu trên các khối 5-s và kiểm tra độ ổn định phương sai trên các khối 10-s theo bài báo chính. Tiêu chí phương sai chỉ đánh giá độ ổn định của phương sai, không xác lập tính dừng đầy đủ của tín hiệu.

\begin{table}[t]
\centering
\caption{Số window được giữ lại theo thời lượng và trạng thái.}\label{tab:supp_retained}
\footnotesize
\setlength{\tabcolsep}{3pt}
\renewcommand{\arraystretch}{1.05}
\begin{tabular*}{\linewidth}{@{\extracolsep{\fill}}lrrr@{}}
\toprule
\makecell{Thời lượng\\(s)} & Awake & Drowsy & Tổng\\
\midrule
30 & 1,192 & 610 & 1,802\\
60 & 596 & 305 & 901\\
120 & 269 & 136 & 405\\
180 & 161 & 77 & 238\\
\bottomrule
\end{tabular*}
\end{table}
\section{Cơ sở lựa chọn tham số nhúng}\label{sec:supp_2}
AMI được ước lượng bằng KSG ($k=3$); độ trễ được chọn tại cực tiểu cục bộ đầu tiên nằm bên trong khoảng khảo sát và có giá trị thấp hơn cả hai điểm liền kề. AMI xác định được độ trễ cho toàn bộ 901/901 window đã xử lý dài 60-s. Trung vị của các trung vị theo session là cơ sở chọn độ trễ chung $\tau=0.16$ s; 14/20 trung vị session nằm trong khoảng $0.16\pm0.04$ s. FNN duy trì thấp quanh $m=7$--9, do đó $m=8$ được chọn làm số chiều nhúng chung. Đây là các lựa chọn thực hành dùng cho cả hai trạng thái, không phải bộ tham số tối ưu duy nhất hay được tối ưu riêng theo trạng thái; các kết quả chẩn đoán này không chứng minh hỗn loạn.

\section{Chi tiết triển khai và hiệu chuẩn Simplex}\label{sec:supp_3}
Thư viện và các đích dự báo đều nằm trong từng window đã chuẩn hóa. Dự báo dùng khoảng cách Euclidean, trọng số láng giềng theo hàm mũ được chuẩn hóa và quy trình leave-one-out. Bảng~\ref{tab:supp_simplex} liệt kê cấu hình cùng toàn bộ horizon dự báo; Mean CC và Mean NRMSE được tính theo định nghĩa trong bài báo chính. Khi hiệu chuẩn khoảng loại trừ theo thời gian, $W=0$--2 s được khảo sát; thay đổi lớn nhất giữa hai mức liền kề xảy ra tại $0.6\to0.8$ s. Chọn $W=1.0$ s vì đây là mức đầu tiên mà kết quả hiệu chuẩn tiếp tục ổn định ở tất cả các mức lớn hơn đã thử. Hiệu chuẩn này mô tả đặc tính của bộ ước lượng, không bổ sung kiểm định về khác biệt giữa hai trạng thái.

\begin{table}[t]
\centering
\caption{Cấu hình Simplex và danh sách đầy đủ các horizon dự báo.}\label{tab:supp_simplex}
\footnotesize
\setlength{\tabcolsep}{3pt}
\renewcommand{\arraystretch}{1.05}
\begin{tabularx}{\linewidth}{@{}lY@{}}
\toprule
Tham số & Cấu hình\\
\midrule
Nhúng & $m=8$, $\tau=0.16$ s\\
Láng giềng & $k=m+1=9$; thư viện trong window\\
Loại trừ & $W=1.0$ s\\
Horizon & \makecell[l]{0.04, 0.08, 0.12, 0.16, 0.20, 0.28,\\0.40, 0.60, 0.80, 1.00, 1.20, 1.60,\\2.00, 2.40, 2.80, 3.20, 3.60, 4.00 s}\\
Chỉ số & Mean CC; Mean NRMSE\\
\bottomrule
\end{tabularx}
\end{table}
\section{Chẩn đoán bộ ước lượng RQA và LLE}\label{sec:supp_4}
RQA sử dụng ngưỡng khoảng cách Euclidean riêng cho từng window để đạt $\mathrm{RR}=0.02$, với khoảng loại trừ $W=(m-1)\tau_{\mathrm{samples}}$ và $l_{\min}=v_{\min}=2$. RR thực tế gần với giá trị mục tiêu. Các chỉ số đường chéo (DET/$L_{\mathrm{mean}}$) và đường thẳng đứng (LAM/TT) mô tả cấu trúc tái diễn; độ dài đường được tính theo chỉ số mẫu. DET không đo tính tất định vật lý của hệ.

LLE theo phương pháp Rosenstein sử dụng khoảng loại trừ dựa trên chu kỳ trung bình phổ, thời gian theo dõi tối đa 5 s và khoảng khớp 0.80--1.30 s. Phép khớp được chấp nhận khi $R^2\geq0.90$, có ít nhất 50 cặp ban đầu, 30 cặp tại mỗi điểm khớp và ba điểm khớp có giá trị hữu hạn. Trong 901 window đầu vào, 872 window đạt yêu cầu. Số cặp ban đầu và số cặp tại điểm khớp nhỏ nhất lần lượt là 1,471 và 1,366; trung vị $R^2$ của phép khớp là 0.973. Các window bị loại do không đạt tiêu chí $R^2$, không phải do thiếu cặp quỹ đạo. LLE ước lượng mức phân kỳ quỹ đạo cục bộ trong một window hữu hạn.

\section{Bằng chứng từ pseudoperiodic surrogates}\label{sec:supp_5}
Mỗi window được so sánh với $M=39$ PPS theo giả thuyết không pseudoperiodic có nhiễu đã kiểm định. Kiểm định thứ hạng hai phía cho mẫu hữu hạn tính các giá trị bằng nhau vào cả hai đuôi. Giá trị $p$ nhỏ nhất có thể đạt là 0.05; tiêu chí bác bỏ là $p\leq0.05$, không hiệu chỉnh ở cấp window. Sáu chỉ số ngoài LLE dùng chung một tập PPS. LLE dùng một tập riêng, với 579 phép khớp gốc Awake/293 Drowsy đã được chấp nhận; các phép khớp LLE trên surrogate không áp dụng tiêu chí $R^2$ của dữ liệu gốc.

Bảng~\ref{tab:supp_pps} phân biệt tỷ lệ bác bỏ gộp với trung vị tỷ lệ bác bỏ theo session, đồng thời giữ các trường hợp bác bỏ ngược hướng. Cả 12 so sánh ở cấp session về chênh lệch giữa giá trị gốc và trung vị PPS (sáu chỉ số ngoài LLE ở cả hai trạng thái) đều có $q_{\mathrm{BH}}<0.05$ sau hiệu chỉnh BH12. Kết quả bác bỏ PPS cho thấy dữ liệu lệch khỏi giả thuyết không pseudoperiodic có nhiễu đã kiểm định. Điều này không đồng nghĩa với việc chứng minh hỗn loạn tất định hay loại trừ mọi mô hình ngẫu nhiên có cấu trúc; tỷ lệ bác bỏ trong từng trạng thái không kiểm định chênh lệch Awake--Drowsy ghép cặp.

\begin{table*}[t]
\centering
\caption{Tóm tắt kết quả bác bỏ PPS của cả bảy chỉ số.}\label{tab:supp_pps}
\footnotesize
\setlength{\tabcolsep}{3pt}
\renewcommand{\arraystretch}{1.05}
\begin{tabular*}{\linewidth}{@{\extracolsep{\fill}}llrrrrr@{}}
\toprule
Chỉ số & Trạng thái & \makecell{Window\\hợp lệ} & \makecell{Bác bỏ\\$n$} & \makecell{Gộp\\(\%)} & \makecell{Trung vị tỷ lệ bác bỏ\\theo session (\%)} & \makecell{Số bác bỏ\\ngược hướng}\\
\midrule
Mean CC & Awake & 596 & 374 & $62.75$ & $60.66$ & 0\\
Mean CC & Drowsy & 305 & 229 & $75.08$ & $71.01$ & 0\\
Mean NRMSE & Awake & 596 & 542 & $90.94$ & $91.59$ & 0\\
Mean NRMSE & Drowsy & 305 & 291 & $95.41$ & $100.00$ & 0\\
DET & Awake & 596 & 595 & $99.83$ & $100.00$ & 0\\
DET & Drowsy & 305 & 305 & $100.00$ & $100.00$ & 0\\
$L_{\mathrm{mean}}$ & Awake & 596 & 594 & $99.66$ & $100.00$ & 0\\
$L_{\mathrm{mean}}$ & Drowsy & 305 & 305 & $100.00$ & $100.00$ & 0\\
LAM & Awake & 596 & 537 & $90.10$ & $95.50$ & 0\\
LAM & Drowsy & 305 & 273 & $89.51$ & $92.38$ & 1\\
TT & Awake & 596 & 580 & $97.32$ & $100.00$ & 0\\
TT & Drowsy & 305 & 295 & $96.72$ & $100.00$ & 0\\
LLE & Awake & 579 & 545 & $94.13$ & $96.30$ & 0\\
LLE & Drowsy & 293 & 287 & $97.95$ & $100.00$ & 0\\
\bottomrule
\end{tabular*}
\par\smallskip
\raggedright Dấu kỳ vọng của chênh lệch giữa giá trị gốc và trung vị PPS: dương với Mean CC, DET, $L_{\mathrm{mean}}$, LLE; âm với Mean NRMSE, LAM, TT. Trung vị tỷ lệ bác bỏ theo session được tính trên 20 session ghi nhận.
\end{table*}
\section{Phân tích tính ổn định của kết quả và độ nhạy}\label{sec:supp_6}
Bảng~\ref{tab:supp_sensitivity} và~\ref{tab:supp_sensitivity_estimators} bổ sung độ bất định và kết quả kiểm định cho các phân tích độ nhạy trên cùng bộ dữ liệu, không phải các lần nghiên cứu lặp lại độc lập. $\Delta$ là trung vị chênh lệch Drowsy trừ Awake ghép cặp giữa các giá trị tổng hợp theo session và trạng thái. CI phân vị bootstrap 95\% và kiểm định Wilcoxon signed-rank tuân theo bài báo chính. Mỗi phân tích chỉ thay đổi một thiết lập; việc giữ nguyên hướng hiệu ứng được xét riêng với độ lớn, độ bất định và bằng chứng thống kê.

\subsection{Độ nhạy theo độ dài window và quy tắc gán nhãn}\label{sec:supp_6_1}
Hướng hiệu ứng trung vị của cả bốn chỉ số chính không đổi qua các window 30/60/120/180 s, nhưng độ lớn và độ bất định thay đổi; CI của DET chứa giá trị không tại 30/180 s (Phần A). Window dài hơn làm giảm số window sẵn có (Bảng~\ref{tab:supp_retained}); 60 s là cấu hình tham chiếu phù hợp cho thực hành, không phải độ dài tối ưu cho mọi trường hợp. P0 là bộ nhãn gốc dùng trong phân tích chính. T30 và T60 loại các mẫu nằm trong khoảng $\pm30$ và $\pm60$ s quanh mỗi thời điểm chuyển Awake--Drowsy hoặc Drowsy--Awake quan sát được; S3 và S5 chỉ giữ các đoạn cùng trạng thái kéo dài ít nhất 180 và 300 s, tương ứng. P0/T30/T60/S3/S5 lần lượt giữ 901/860/787/900/846 window sau QC chung, với $n=20$ session ghép cặp, ngoại trừ S5 ($n=19$). Các quy tắc thay thế giữ nguyên hướng của cả bốn hiệu ứng và đều đạt ý nghĩa thống kê sau BH4 (Phần B), dù CI của Mean CC theo T60 và LLE theo T30 chứa giá trị không. CI phân vị của hiệu ứng trung vị và kết quả kiểm định dựa trên thứ hạng không nhất thiết trùng khớp hoàn toàn.

\subsection{Độ nhạy theo tham số nhúng}\label{sec:supp_6_2}
Mean CC, Mean NRMSE và LLE giữ nguyên hướng hiệu ứng qua các độ trễ và số chiều đã thử (Phần C). DET đổi sang hiệu ứng dương gần không tại $\tau=0.12$ s ($+0.0025$) và 0.20 s ($+0.0015$); cả hai CI đều chứa giá trị không. DET vẫn có hiệu ứng âm trong khoảng $m=7$--9, nhưng không đạt ý nghĩa thống kê sau BH4 tại $m=7$ và 9. Vì vậy, DET nhạy với độ trễ nhúng; hướng hiệu ứng không đổi theo số chiều không có nghĩa là độ lớn bằng nhau hay bằng chứng thống kê đồng đều.

\subsection{Độ nhạy theo thiết lập bộ ước lượng}\label{sec:supp_6_3}
DET vẫn có hiệu ứng âm ở các thiết lập RR/Theiler đã thử, nhưng bằng chứng thống kê yếu hơn tại RR=0.01 (Phần D). LAM/TT không đạt ý nghĩa thống kê ở các thiết lập này. LLE có hiệu ứng âm ở tất cả các khoảng khớp, ngưỡng $R^2$ và hệ số Theiler đã thử; mọi CI của LLE đều không chứa giá trị không và $p<0.05$ khi chưa hiệu chỉnh. Thiết lập khớp/QC làm thay đổi độ lớn hiệu ứng và số window được chấp nhận. Các so sánh bộ ước lượng chỉ báo cáo p chưa hiệu chỉnh. Hệ số Theiler được nhân với khoảng loại trừ danh định của từng bộ ước lượng.

\begin{table*}[!t]
% Window/label and reconstruction/estimator panels share Table S4.
\renewcommand{\thetable}{S\arabic{table}A}
\centering
\caption{Độ nhạy theo độ dài window và quy tắc nhãn.}\label{tab:supp_sensitivity}
\footnotesize
\setlength{\tabcolsep}{1.5pt}
\renewcommand{\arraystretch}{1.12}
\begin{minipage}[t]{0.49\textwidth}
\raggedright\textbf{A. Độ dài window}\par\smallskip
\begin{tabular*}{\linewidth}{@{\extracolsep{\fill}}llrcr@{}}
\toprule
Chỉ số & s & $\Delta$ & 95\% CI & $p$ thô\\
\midrule
CC & 30 & $-0.0147$ & $[-0.0712,\,-0.0039]$ & $0.03276825$\\
CC & 60 & $-0.0339$ & $[-0.0595,\,-0.0043]$ & $0.02148438$\\
CC & 120 & $-0.0271$ & $[-0.0490,\,-0.0156]$ & $0.01361656$\\
CC & 180 & $-0.0428$ & $[-0.0516,\,-0.0133]$ & $0.00943565$\\
NRMSE & 30 & $+0.0203$ & $[+0.0009,\,+0.0554]$ & $0.00943565$\\
NRMSE & 60 & $+0.0294$ & $[0.0047,\,0.0530]$ & $0.00729561$\\
NRMSE & 120 & $+0.0321$ & $[+0.0196,\,+0.0495]$ & $0.00830841$\\
NRMSE & 180 & $+0.0321$ & $[+0.0186,\,+0.0465]$ & $0.00729561$\\
DET & 30 & $-0.0169$ & $[-0.0318,\,+0.0002]$ & $0.01531219$\\
DET & 60 & $-0.0186$ & $[-0.0364,\,-0.0069]$ & $0.02148438$\\
DET & 120 & $-0.0201$ & $[-0.0399,\,-0.0134]$ & $0.00638962$\\
DET & 180 & $-0.0248$ & $[-0.0416,\,+0.0008]$ & $0.05316925$\\
LLE & 30 & $-0.0325$ & $[-0.0549,\,-0.0066]$ & $0.00315285$\\
LLE & 60 & $-0.0376$ & $[-0.0537,\,-0.0230]$ & $0.00070763$\\
LLE & 120 & $-0.0442$ & $[-0.0591,\,-0.0216]$ & $0.00232506$\\
LLE & 180 & $-0.0498$ & $[-0.0584,\,-0.0213]$ & $0.00058556$\\
\bottomrule
\end{tabular*}
\end{minipage}\hfill\begin{minipage}[t]{0.49\textwidth}
\raggedright\textbf{B. Quy tắc nhãn}\par\smallskip
\begin{tabular*}{\linewidth}{@{\extracolsep{\fill}}llrcr@{}}
\toprule
Quy tắc & Chỉ số & $\Delta$ & 95\% CI & $q_{\mathrm{BH}}$\\
\midrule
\textbf{P0} & CC & $-0.0339$ & $[-0.0595,\,-0.0043]$ & $0.0376$\\
\textbf{P0} & NRMSE & $+0.0294$ & $[0.0047,\,0.0530]$ & $0.0255$\\
\textbf{P0} & DET & $-0.0186$ & $[-0.0364,\,-0.0069]$ & $0.0376$\\
\textbf{P0} & LLE & $-0.0376$ & $[-0.0537,\,-0.0230]$ & $0.0050$\\
T30 & CC & $-0.0482$ & $[-0.0689,\,-0.0152]$ & $0.0097$\\
T30 & NRMSE & $+0.0406$ & $[+0.0190,\,+0.0570]$ & $0.0097$\\
T30 & DET & $-0.0183$ & $[-0.0332,\,-0.0011]$ & $0.0111$\\
T30 & LLE & $-0.0322$ & $[-0.0535,\,+0.0020]$ & $0.0215$\\
T60 & CC & $-0.0251$ & $[-0.0561,\,+0.0059]$ & $0.0400$\\
T60 & NRMSE & $+0.0278$ & $[+0.0027,\,+0.0479]$ & $0.0256$\\
T60 & DET & $-0.0167$ & $[-0.0294,\,-0.0066]$ & $0.0256$\\
T60 & LLE & $-0.0393$ & $[-0.0636,\,-0.0075]$ & $0.0093$\\
S3 & CC & $-0.0361$ & $[-0.0548,\,-0.0033]$ & $0.0136$\\
S3 & NRMSE & $+0.0357$ & $[+0.0032,\,+0.0457]$ & $0.0128$\\
S3 & DET & $-0.0189$ & $[-0.0326,\,-0.0069]$ & $0.0136$\\
S3 & LLE & $-0.0344$ & $[-0.0495,\,-0.0104]$ & $0.0041$\\
S5 & CC & $-0.0337$ & $[-0.0495,\,-0.0077]$ & $0.0141$\\
S5 & NRMSE & $+0.0367$ & $[+0.0056,\,+0.0504]$ & $0.0092$\\
S5 & DET & $-0.0213$ & $[-0.0305,\,-0.0049]$ & $0.0141$\\
S5 & LLE & $-0.0321$ & $[-0.0499,\,-0.0147]$ & $0.0056$\\
\bottomrule
\end{tabular*}
\end{minipage}
\par\smallskip
\raggedright CC/NRMSE là Mean CC/Mean NRMSE. Hiệu ứng và CI của LLE có đơn vị s$^{-1}$; các hiệu ứng chính còn lại không có đơn vị. Các dòng P0 in đậm là giá trị tham chiếu chính của bài báo, dùng BH7; các quy tắc nhãn thay thế dùng BH4 riêng trong từng quy tắc. Phần A chỉ báo cáo p chưa hiệu chỉnh.

\end{table*}

\begin{table*}[!t]
\ContinuedFloat
\renewcommand{\thetable}{S\arabic{table}B}
\centering
\caption{Độ nhạy theo tham số nhúng và thiết lập bộ ước lượng.}\label{tab:supp_sensitivity_estimators}
\footnotesize
\setlength{\tabcolsep}{1.5pt}
\renewcommand{\arraystretch}{1.12}
\begin{minipage}[t]{0.49\textwidth}
\raggedright\textbf{C. Độ trễ ($m=8$)}\par\smallskip
\begin{tabular*}{\linewidth}{@{\extracolsep{\fill}}llrcr@{}}
\toprule
$\tau$ (s) & Chỉ số & $\Delta$ & 95\% CI & $q_{\mathrm{BH}}$\\
\midrule
$0.12$ & CC & $-0.0317$ & $[-0.0522,\,+0.0094]$ & $0.048312$\\
$0.12$ & NRMSE & $+0.0321$ & $[-0.0042,\,+0.0492]$ & $0.048312$\\
$0.12$ & DET & $+0.0025$ & $[-0.0078,\,+0.0076]$ & $0.784126$\\
$0.12$ & LLE & $-0.0361$ & $[-0.0567,\,+0.0042]$ & $0.048312$\\
\textbf{0.16} & CC & $-0.0339$ & $[-0.0595,\,-0.0043]$ & $0.0376$\\
\textbf{0.16} & NRMSE & $+0.0294$ & $[0.0047,\,0.0530]$ & $0.0255$\\
\textbf{0.16} & DET & $-0.0186$ & $[-0.0364,\,-0.0069]$ & $0.0376$\\
\textbf{0.16} & LLE & $-0.0376$ & $[-0.0537,\,-0.0230]$ & $0.0050$\\
$0.20$ & CC & $-0.0316$ & $[-0.0529,\,-0.0008]$ & $0.020416$\\
$0.20$ & NRMSE & $+0.0296$ & $[+0.0056,\,+0.0476]$ & $0.011162$\\
$0.20$ & DET & $+0.0015$ & $[-0.0099,\,+0.0148]$ & $0.985435$\\
$0.20$ & LLE & $-0.0464$ & $[-0.0572,\,-0.0240]$ & $0.000145$\\
\bottomrule
\end{tabular*}
\end{minipage}\hfill\begin{minipage}[t]{0.49\textwidth}
\raggedright\textbf{C. Số chiều ($\tau=0.16$ s)}\par\smallskip
\begin{tabular*}{\linewidth}{@{\extracolsep{\fill}}llrcr@{}}
\toprule
$m$ & Chỉ số & $\Delta$ & 95\% CI & $q_{\mathrm{BH}}$\\
\midrule
7 & CC & $-0.0307$ & $[-0.0565,\,+0.0045]$ & $0.035522$\\
7 & NRMSE & $+0.0377$ & $[-0.0023,\,+0.0561]$ & $0.035522$\\
7 & DET & $-0.0145$ & $[-0.0303,\,-0.0047]$ & $0.053169$\\
7 & LLE & $-0.0285$ & $[-0.0549,\,-0.0076]$ & $0.019440$\\
9 & CC & $-0.0251$ & $[-0.0492,\,+0.0055]$ & $0.039434$\\
9 & NRMSE & $+0.0315$ & $[+0.0043,\,+0.0483]$ & $0.016617$\\
9 & DET & $-0.0200$ & $[-0.0321,\,-0.0052]$ & $0.058258$\\
9 & LLE & $-0.0247$ & $[-0.0534,\,-0.0102]$ & $0.002342$\\
\bottomrule
\end{tabular*}
\end{minipage}
\par\medskip
\begin{minipage}[t]{0.49\textwidth}
\raggedright\textbf{D. RQA: DET}\par\smallskip
\begin{tabular*}{\linewidth}{@{\extracolsep{\fill}}lrcr@{}}
\toprule
Cấu hình & $\Delta$ & 95\% CI & $p$ thô\\
\midrule
RR 0.01 & $-0.015214$ & $[-0.033764,\,+0.002064]$ & $0.082550$\\
Danh định & $-0.0186$ & $[-0.0364,\,-0.0069]$ & $0.021484$\\
RR 0.03 & $-0.017813$ & $[-0.033859,\,-0.007863]$ & $0.021484$\\
Theiler 0.75 & $-0.019191$ & $[-0.036836,\,-0.006809]$ & $0.021484$\\
Theiler 1.25 & $-0.019033$ & $[-0.037873,\,-0.007730]$ & $0.017181$\\
\bottomrule
\end{tabular*}
\end{minipage}\hfill\begin{minipage}[t]{0.49\textwidth}
\raggedright\textbf{D. LLE}\par\smallskip
\begin{tabular*}{\linewidth}{@{\extracolsep{\fill}}lrcrr@{}}
\toprule
Cấu hình & $\Delta$ & 95\% CI & $p$ thô & Đạt/tổng\\
\midrule
Khớp 0.60--1.10 & $-0.034980$ & $[-0.064607,\,-0.019121]$ & $0.000261$ & 802/901\\
Danh định & $-0.0376$ & $[-0.0537,\,-0.0230]$ & $0.000708$ & 872/901\\
Khớp 1.00--1.50 & $-0.030357$ & $[-0.045778,\,-0.003941]$ & $0.008308$ & 889/901\\
$R^2$ 0.95 & $-0.030618$ & $[-0.041701,\,-0.013987]$ & $0.002325$ & 734/901\\
Theiler 0.75 & $-0.037570$ & $[-0.053698,\,-0.022976]$ & $0.000708$ & 872/901\\
Theiler 1.25 & $-0.029763$ & $[-0.056095,\,-0.016331]$ & $0.002325$ & 874/901\\
\bottomrule
\end{tabular*}
\end{minipage}
\par\smallskip
\raggedright CC/NRMSE là Mean CC/Mean NRMSE. Hiệu ứng và CI của LLE có đơn vị s$^{-1}$; các hiệu ứng chính còn lại không có đơn vị. Các tham số nhúng danh định in đậm là giá trị tham chiếu chính của bài báo, dùng BH7; các thiết lập nhúng thay thế dùng BH4 riêng trong từng thiết lập. Phần D chỉ báo cáo p chưa hiệu chỉnh.
\par\smallskip
\raggedright Phần C: $m=8$ dùng giá trị tham chiếu $\tau=0.16$ ở bảng bên trái. Phần D: RQA danh định dùng RR=0.02 và hệ số Theiler 1.00; LLE danh định dùng khoảng khớp 0.80--1.30 s, $R^2\geq0.90$ và hệ số 1.00. Khoảng khớp tính bằng giây. Mọi thiết lập RQA giữ 901 window; LLE ghi số window đạt yêu cầu trên tổng số đầu vào. Các tham số không được thay đổi giữ ở giá trị danh định.
\end{table*}
\FloatBarrier
\section{Phụ thuộc giữa các session khi chưa biết liên kết người tham gia}\label{sec:supp_7}
Không có thông tin liên kết từng session với người tham gia. Các cách ghép cặp hoàn hảo giả định chia 20 session thành 10 cụm, mỗi cụm gồm hai session: 100,000 lượt lấy mẫu tạo ra 99,990 cách ghép khác nhau. Chênh lệch theo session được lấy trung bình trong từng cụm rồi kiểm định trên các cụm, với hiệu chỉnh BH7 trong mỗi cách ghép. Hướng hiệu ứng trung vị của cả bốn chỉ số chính đều không đổi trong mọi cách ghép giả định đã lấy mẫu (Bảng~\ref{tab:supp_dependence}), nhưng kết quả đạt ý nghĩa thống kê ít nhất quán hơn so với phân tích ở cấp session. Ước lượng tại 50k và 100k lượt rất gần nhau, cho thấy độ ổn định Monte Carlo. Các tần suất này mô tả thiết kế ghép cặp giả định, không phải xác suất kết quả thực có ý nghĩa thống kê. Cách nhóm giả định không khôi phục danh tính người tham gia và không cho phép suy luận ở cấp người tham gia.

\begin{table}[!t]
\centering
\caption{Các hiệu ứng chính khi nhóm session theo cách ghép giả định.}\label{tab:supp_dependence}
\footnotesize
\setlength{\tabcolsep}{3pt}
\renewcommand{\arraystretch}{1.05}
\begin{tabular}{@{}lrr@{}}
\toprule
Chỉ số & \makecell{Dấu không đổi\\(\%)} & \makecell{$q_{\mathrm{BH}}<0.05$\\tần suất (\%)}\\
\midrule
Mean CC & $100.00$ & $19.36$\\
Mean NRMSE & $100.00$ & $27.82$\\
DET & $100.00$ & $22.59$\\
LLE & $100.00$ & $79.69$\\
\bottomrule
\end{tabular}
\end{table}
\FloatBarrier
\section{Phân tích động lực học ký hiệu mang tính thăm dò}\label{sec:supp_8}
Khoảng thời gian giữa các đỉnh liên tiếp (PPI) được trích từ PPG đã xử lý và lượng tử hóa trong từng window được chấp nhận thành sáu mức có độ rộng bằng nhau, từ giá trị nhỏ nhất đến lớn nhất. Các bộ ba ký hiệu chồng lấp được phân vào 0V, 1V hoặc 2V khi không có, có một hoặc có cả hai chênh lệch giữa các ký hiệu liền kề khác không; 2V bao gồm các thay đổi cùng dấu và ngược dấu. Tỷ lệ xuất hiện được tính cho từng window, sau đó lấy trung vị trong mỗi session và trạng thái.

Bảng~\ref{tab:supp_symbolic} trình bày trung vị chênh lệch ghép cặp theo điểm phần trăm. Cả chín kiểm định nhóm ký hiệu--thời lượng đều không đạt ý nghĩa thống kê sau BH3 trong từng thời lượng. Các kết quả thăm dò này nằm ngoài kết luận chính dựa trên bốn chỉ số; tỷ lệ ký hiệu không đo trực tiếp hoạt động giao cảm hay phế vị.

\begin{table}[!t]
\centering
\caption{Trung vị hiệu ứng ký hiệu trong phân tích thăm dò (điểm phần trăm).}\label{tab:supp_symbolic}
\footnotesize
\setlength{\tabcolsep}{3pt}
\renewcommand{\arraystretch}{1.05}
\begin{tabular}{@{}lrrr@{}}
\toprule
Thời lượng (s) & 0V $\Delta$ & 1V $\Delta$ & 2V $\Delta$\\
\midrule
60 & $+4.04$ & $+0.41$ & $-0.91$\\
120 & $+2.32$ & $+0.03$ & $-2.30$\\
180 & $+4.58$ & $-0.06$ & $-2.67$\\
\bottomrule
\end{tabular}
\end{table}
\FloatBarrier
\end{document}
